\section{Experiments}
\label{sec:experiments}

\subsection{Setup}

\textbf{Dataset.} We evaluate on Zhuhai, China, using a POI database of 2129 entries across categories including scenic spots, restaurants, cultural sites, natural areas, and entertainment venues. Each POI includes geographic coordinates, operating hours, ratings, price ranges, and 6-dimensional emotion tags (excitement, tranquility, sociability, culture depth, surprise, physical demand).

We use three evaluation scales: (1)~5 core scenarios for detailed per-dimension analysis with LLM judge; (2)~30 scenarios for development-stage optimization tracking; and (3)~100 scenarios for final large-scale evaluation.

\textbf{5 Core Scenarios.} Each represents one of the five scenario types:
\begin{itemize}[leftmargin=*]
    \item S1: ``Couple's one-day Zhuhai tour, budget 500 yuan, like photo spots'' (Sightseeing)
    \item S2: ``Take 6-year-old to Chimelong Ocean Kingdom, budget 1000 yuan'' (Destination)
    \item S3: ``Zhuhai food one-day tour, want seafood and local specialties'' (Food)
    \item S4: ``Check in all famous Zhuhai spots in one day, tight schedule'' (Speed-Run)
    \item S5: ``Two-day Zhuhai tour, slow pace, like parks and seaside'' (Leisure)
\end{itemize}

\textbf{30-Scenario Baseline.} We also maintain a 30-scenario stress test covering all five scenario types (6 scenarios each), used for development-stage optimization tracking. The baseline (commit \texttt{c9dad35}) achieved 25/30 pass rate with an average score of 7.3, with 5 failing scenarios: friend gathering (\#11), seafood feast (\#21), food exploration (\#22), budget food tour (\#25), and vague request (\#30). The weakest dimension across all scenarios was rhythm (6.5 average).

\textbf{100-Scenario Final Evaluation.} The final evaluation uses 100 scenarios (20 per scenario type), generated from structured templates varying user profile, budget, time constraints, and preference dimensions. Each scenario is scored on a 0--100 scale with grade thresholds: S ($\geq$90), A (80--89), B (70--79), C (60--69), D (50--59), F ($<$50). Scoring dimensions include POI count, POI quality, geographic continuity, diversity, completeness, time feasibility, and category match.

\textbf{Evaluation.} We employ a dual evaluation protocol:
\begin{enumerate}[leftmargin=*]
    \item \textbf{LLM Judge}: An independent LLM evaluator (DeepSeek-V4-Flash via API) with scenario-aware scoring rubrics scores each route on four dimensions (0--10 scale): intent match, POI quality, geographic continuity, and scene diversity. The overall score is the unweighted average. We run each evaluation 2 times and report the lower-score run (conservative estimate) to reduce LLM scoring variance~\cite{gu2024llmjudge}.
    \item \textbf{Algorithmic Evaluation}: A rule-based evaluation framework computes category coverage (30\%), emotion fit (30\%), feasibility (20\%), and intent match (20\%) without LLM dependency, providing a deterministic baseline for reproducibility.
\end{enumerate}

\subsection{Main Results}

Table~\ref{tab:scenario_changes} shows the per-scenario scores before and after each optimization. Table~\ref{tab:main_results} summarizes the final 100-scenario evaluation across all dimensions.

\begin{table}[!t]
\centering
\caption{Per-scenario evaluation scores (0--10) for CityFlow on 5 core scenarios. Scores from LLM judge (Qwen3.5-35B-A3B). Test run: 2026-06-02.}
\label{tab:main_results}
\resizebox{\columnwidth}{!}{\begin{tabular}{@{}lcccccc@{}}
\toprule
\textbf{Scenario} & \textbf{Intent} & \textbf{POI} & \textbf{Geo} & \textbf{Diversity} & \textbf{Overall} & \textbf{Stops} \\
\midrule
S1: Couple Tour & 8 & 7 & 6 & 7 & \textbf{7} & 6 \\
S2: Ocean Kingdom & 9 & 8 & 9 & 5 & \textbf{8} & 3 \\
S3: Food Tour & 8 & 7 & 6 & 5 & \textbf{7} & 6 \\
S4: Speed-Run & 8 & 6 & 5 & 5 & \textbf{7} & 5 \\
S5: Leisure & 9 & 8 & 7 & 6 & \textbf{8} & 4 \\
\midrule
\textbf{Average} & \textbf{8.4} & \textbf{7.2} & \textbf{6.6} & \textbf{5.6} & \textbf{7.4} & \textbf{4.8} \\
\bottomrule
\end{tabular}}
\end{table}

CityFlow achieves an average overall score of 7.4/10 with 5/5 pass rate (all scenarios $\geq$ 6.5). Intent match (8.4) is the strongest dimension, while scene diversity (5.6) is the weakest, particularly for destination-focused and food scenarios where fewer POI types are appropriate. Geographic continuity (6.6) shows room for improvement, with speed-run and food scenarios occasionally exhibiting backtracking.

\subsection{100-Scenario Large-Scale Evaluation}

Table~\ref{tab:100_overview} reports the overall results of the 100-scenario evaluation.

\begin{table}[!t]
\centering
\caption{100-scenario evaluation summary. Test: 2026-06-02, 4 workers, total 2263.6s.}
\label{tab:100_overview}
\resizebox{\columnwidth}{!}{\begin{tabular}{@{}lc@{}}
\toprule
\textbf{Metric} & \textbf{Value} \\
\midrule
Total scenarios & 100 \\
Route generation success & 90/100 (10 failed) \\
Average score (0--100) & 74.3 \\
Average latency per scenario & 70.3s \\
\midrule
Grade distribution & S:12, A:24, B:23, C:30, D:10, F:1 \\
\bottomrule
\end{tabular}}
\end{table}

Table~\ref{tab:100_by_type} breaks down results by scenario type.

\begin{table}[!t]
\centering
\caption{100-scenario results by scenario type (20 scenarios each).}
\label{tab:100_by_type}
\resizebox{\columnwidth}{!}{\begin{tabular}{@{}lccc@{}}
\toprule
\textbf{Type} & \textbf{Avg Score} & \textbf{Pass} & \textbf{S/A/B/C/D/F} \\
\midrule
Food & 76.8 & 19/20 & 4/3/6/6/1/0 \\
Destination & 68.3 & 15/20 & 0/3/5/7/5/0 \\
Speed-Run & 76.5 & 20/20 & 3/5/5/7/0/0 \\
Leisure & 73.5 & 17/20 & 1/9/2/5/2/1 \\
Sightseeing & 76.4 & 18/20 & 4/4/5/5/2/0 \\
\bottomrule
\end{tabular}}
\end{table}

Key observations:
\begin{itemize}[leftmargin=*]
    \item \textbf{Speed-Run achieves 100\% pass rate} (20/20), validating the time-smoothing and refusal mechanisms that were specifically designed for this scenario type.
    \item \textbf{Destination is the weakest type} (68.3 avg, 15/20 pass, 5 D-grades). This is expected: destination scenarios involve single-venue deep visits (e.g., hot springs, museums) where the expert-ensemble architecture has less room to differentiate, and the POI database has fewer entries for niche attractions.
    \item \textbf{Food, Sightseeing, and Speed-Run} all average $>$76, demonstrating strong performance across the most common route planning use cases.
    \item \textbf{The single F-grade} (29.2, scenario \#61: ``Couple's weekend leisure, slow stroll, relax'') failed because the system produced only 2 stops for a vague leisure request, resulting in low completeness and diversity scores.
\end{itemize}

Table~\ref{tab:100_dims} reports per-dimension scores by scenario type, revealing systematic strengths and weaknesses.

\begin{table}[!t]
\centering
\caption{Per-dimension average scores by scenario type (0--100 scale).}
\label{tab:100_dims}
\resizebox{\columnwidth}{!}{\begin{tabular}{@{}lccccccc@{}}
\toprule
\textbf{Type} & \textbf{POI\#} & \textbf{POI-Q} & \textbf{Geo} & \textbf{Div} & \textbf{Comp} & \textbf{Time} & \textbf{Cat} \\
\midrule
Food & 13 & 82 & 62 & 57 & 69 & 40 & 85 \\
Destination & 4 & 85 & 62 & 65 & 67 & 63 & 77 \\
Speed-Run & 4 & 83 & 49 & 73 & 78 & 44 & 68 \\
Leisure & 32 & 84 & 81 & 65 & 69 & 72 & 88 \\
Sightseeing & 9 & 88 & 77 & 69 & 75 & 68 & 81 \\
\midrule
\textbf{Average} & \textbf{12} & \textbf{84} & \textbf{66} & \textbf{66} & \textbf{72} & \textbf{57} & \textbf{80} \\
\bottomrule
\end{tabular}}
\end{table}

\textbf{POI Quality} (avg 84) and \textbf{Category Match} (avg 80) are consistently strong, confirming that the expert-ensemble architecture effectively selects relevant POIs and matches user intent. \textbf{Time Feasibility} (avg 57) is the weakest dimension across all types---LLM-generated time allocations remain the primary source of quality loss, despite algorithmic smoothing. \textbf{Geographic Continuity} (avg 66) shows improvement potential, particularly for Speed-Run routes where tight schedules amplify spatial inefficiency.

\textbf{Leisure} scenarios have the highest stop count (avg 32) and best geographic continuity (81) and time feasibility (72), because leisure routes naturally cluster around walkable areas with relaxed time constraints. Conversely, \textbf{Speed-Run} routes have the lowest geographic continuity (49) due to aggressive multi-stop scheduling across dispersed locations.

\subsection{Error Analysis}

We analyze the 10 failed scenarios (route generation errors) and 8 D-grade scenarios to identify systematic failure patterns. Table~\ref{tab:error_analysis} categorizes the failures.

\begin{table}[!t]
\centering
\caption{Failure mode analysis. 10 route generation failures + 8 D-grade scenarios from 100-scenario evaluation.}
\label{tab:error_analysis}
\resizebox{\columnwidth}{!}{\begin{tabular}{@{}llc@{}}
\toprule
\textbf{Failure Mode} & \textbf{Count} & \textbf{Example} \\
\midrule
Over-aggressive time pruning & 7 & \#7: ``only 2 hours'' $\rightarrow$ 1 stop removed \\
Insufficient stops for vague intent & 2 & \#61: ``slow stroll'' $\rightarrow$ only 1 stop (F) \\
Geographic incoherence (islands) & 3 & \#36: Wai Lingding Island $\rightarrow$ 15km gap \\
Category mismatch (niche venues) & 3 & \#33: Watch museum $\rightarrow$ low category match (20) \\
Time feasibility (spas/islands) & 4 & \#39: Hot spring $\rightarrow$ transit time underestimated \\
\bottomrule
\end{tabular}}
\end{table}

\textbf{Over-aggressive time pruning} is the dominant failure mode (7/18 failures). The \texttt{\_smooth\_times()} function correctly identifies that LLM-generated schedules exceed operating hours, but the remedy---removing the offending stop---is too aggressive. For example, scenario \#27 (``evening fireworks at Chimelong'') had 4 stops pruned to 2 because the fireworks show ends at 22:00 and subsequent stops were scheduled after closing. A better approach would be to re-schedule rather than remove.

\textbf{Insufficient stops for vague intent} affects leisure and slow-paced scenarios. Scenario \#61 (``couple weekend leisure, slow stroll'') produced only 1 stop because the system could not decompose a vague request into concrete destinations. The POI quality score was 4.3/100, indicating the single stop was irrelevant. This suggests that very vague requests may benefit from a clarification dialogue rather than direct generation.

\textbf{Geographic incoherence} affects island-based scenarios where inter-island ferry transit is not modeled in the distance calculations. Wai Lingding Island scenarios (\#28, \#36) and Hengqin Island scenarios show large geographic gaps because the haversine distance calculation treats water as traversable by walking.

\subsection{30-Scenario Baseline Progression}

Table~\ref{tab:comparison} documents the progression from the initial solver-based system to the final CityFlow architecture, measured on the 30-scenario stress test at three development milestones.

\begin{table}[!t]
\centering
\caption{System progression on 30 scenarios across three development milestones. Scores are rule-based composite (0--100 scale).}
\label{tab:comparison}
\resizebox{\columnwidth}{!}{\begin{tabular}{@{}llccc@{}}
\toprule
\textbf{Commit} & \textbf{Architecture} & \textbf{Pass Rate} & \textbf{Avg Score} & \textbf{F-Grade} \\
\midrule
Pre-architecture & Solver-based (rule) & 18\% (5/30) & 55.0 & --- \\
\texttt{c9dad35} & Expert-ensemble (pre-optimization) & 83\% (25/30) & 73.0 & 5 \\
\texttt{821b15a} & Expert-ensemble (post-optimization) & 100\% (30/30) & 81.2 & 0 \\
\bottomrule
\end{tabular}}
\end{table}

The progression shows three distinct phases. The solver-based system relied on rule-based POI ranking with TSP-style nearest-neighbor ordering, achieving only 18\% pass rate. The introduction of LLM-based expert-ensemble (commit \texttt{c9dad35}) dramatically improved pass rate to 83\% and average score from 55.0 to 73.0. After seven rounds of targeted optimization (documented in Section~\ref{sec:journey}), the final system achieved 100\% pass rate with 81.2 average score.

The 5 failing scenarios in the pre-optimization baseline reveal specific weaknesses: friend gathering (\#11), seafood feast (\#21), food exploration (\#22), budget food tour (\#25), and vague request (\#30). These were resolved through a combination of prompt engineering (refusal mechanism), data enrichment (52 food POIs), and algorithmic post-processing (time smoothing).

\subsection{Ablation Study}

To understand the contribution of each architectural component, we analyze the impact of key changes documented in our optimization journal (commit \texttt{821b15a}) and optimization log (commit \texttt{d323625}). Table~\ref{tab:ablation} summarizes the observed score changes when components were removed or modified.

\begin{table}[!t]
\centering
\caption{Ablation analysis derived from optimization history. Each row corresponds to a documented change.}
\label{tab:ablation}
\resizebox{\columnwidth}{!}{\begin{tabular}{@{}llcc@{}}
\toprule
\textbf{Change} & \textbf{Source} & \textbf{Pass} & \textbf{Overall} \\
\midrule
Full CityFlow (current) & --- & 5/5 & 7.4 \\
\midrule
w/o scene-aware eval (\texttt{2cf655d} baseline) & journal & 5/5 & 6.8 \\
w/o scene-aware dispatch (\texttt{7edcb97}) & journal & 4/5 & 7.0 \\
w/o time smoothing (\texttt{3183bea} pre-fix) & journal & 3/5 & --- \\
w/o review-rework (\texttt{d323625} \#7 diversity) & opt log & 3/5 & 6.6 \\
Expert rule-ification (\texttt{d323625} \#3) & opt log & 3/5 & 6.4 \\
\bottomrule
\end{tabular}}
\end{table}

Key findings:
\begin{itemize}[leftmargin=*]
    \item \textbf{Time smoothing} has the largest impact. The speed-run scenario scored F~(42) before time smoothing was introduced and A~(85) after (commit \texttt{487f99a}).
    \item \textbf{Scene-aware dispatch} improved pass rate from 2/5 to 4/5 at commit \texttt{7edcb97} by differentiating agent strategies per scenario type.
    \item \textbf{Expert rule-ification} (replacing LLM experts with rules) caused regression from 4/5 to 3/5 (commit \texttt{d323625}, attempt \#3), confirming that LLM-based experts outperform rule-based alternatives even when using inexpensive models.
    \item \textbf{Scene-aware evaluation} (\texttt{2cf655d}) raised the overall score from 6.8 to 7.4 without changing the generation pipeline, demonstrating that evaluation methodology significantly affects reported scores.
\end{itemize}

\subsection{Model Efficiency Analysis}

A key finding of this work is that \textbf{inexpensive small language models can match expensive frontier models when embedded in a well-designed multi-agent architecture}. We conducted a multi-model A/B test (commit \texttt{edfe2c0}) comparing four models as the expert backbone, with all other architectural components held constant. Table~\ref{tab:model_compare} summarizes the results.

\begin{table}[!t]
\centering
\caption{Model comparison. Pass rate and overall score from multi-model A/B test (commit \texttt{edfe2c0}); Qwen3.5-35B-A3B row from 2026-06-02 test. Pricing as of 2025-05.}
\label{tab:model_compare}
\resizebox{\columnwidth}{!}{\begin{tabular}{@{}lccp{2.5cm}@{}}
\toprule
\textbf{Model} & \textbf{Pass} & \textbf{Overall} & \textbf{Price} \\
\midrule
DeepSeek-V4-Flash & 4/5 & 6.4 & \$\textasciitilde\$1/M tokens \\
qwen-turbo & 3/5 & 6.4 & \$\textasciitilde\$0.5/M tokens \\
qwen-flash & 2/5 & 6.4 & \$\textasciitilde\$0.1/M tokens \\
\textbf{qwen3.5-flash} & \textbf{5/5} & \textbf{7.0} & \textbf{\$\textasciitilde\$0.2/M tokens} \\
\midrule
Qwen3.5-35B-A3B (iFlytek) & \textbf{5/5} & \textbf{7.4} & \textbf{\$\textasciitilde\$0.15/M tokens} \\
\bottomrule
\end{tabular}}
\end{table}

The results are striking: qwen3.5-flash (a \$0.2/M model) achieves a higher overall score than DeepSeek-V4-Flash (a \$1/M model), and the Qwen3.5-35B-A3B model---a 35B-parameter MoE model (3B active) deployed via iFlytek Spark API---achieves the best overall score of 7.4 while being approximately 6$\times$ cheaper than DeepSeek.

\textbf{Why does the small model outperform?} We attribute this to three factors:

\begin{enumerate}[leftmargin=*]
    \item \textbf{Task decomposition reduces per-call complexity.} Each expert handles a narrow domain (e.g., ``rank restaurants by sub-category''), making the reasoning task tractable for smaller models. A single LLM must simultaneously handle POI selection, geographic planning, time scheduling, and diversity balancing---a task where frontier models have an advantage.

    \item \textbf{Algorithmic post-processing compensates for reasoning gaps.} The \texttt{\_smooth\_times()} function corrects temporal infeasibility, and the geographic threshold checks catch spatial errors. These deterministic safeguards prevent small-model hallucinations from reaching the final output.

    \item \textbf{Scene-aware prompts reduce ambiguity.} By classifying the user intent first and dispatching domain-specific prompts, each expert receives a focused, unambiguous instruction. This mitigates the instruction-following weakness of smaller models.
\end{enumerate}

This finding has practical implications: a CityFlow deployment using Qwen3.5-35B-A3B processes a single route request in ${\sim}35$s with an estimated cost of \$0.001 per request, compared to ${\sim}48$s and \$0.006 per request with DeepSeek-V4-Flash (commit \texttt{a576cc5}). The expert-ensemble architecture thus enables cost-effective deployment without quality sacrifice.

\subsection{Architecture Evolution}

CityFlow underwent five major architectural iterations before reaching its current form. Table~\ref{tab:arch_evolution} summarizes the evolution.

\begin{table}[!t]
\centering
\caption{Architecture evolution milestones. Each iteration is identified by its commit hash.}
\label{tab:arch_evolution}
\resizebox{\columnwidth}{!}{\begin{tabular}{@{}llll@{}}
\toprule
\textbf{Commit} & \textbf{Architecture} & \textbf{Pass Rate} & \textbf{Overall} \\
\midrule
\texttt{107d75d} & C-version: 7 Agents + LLM coordinator & 4/5 & 7.2 \\
\texttt{7edcb97} & Scene-aware: 5 scene types & 4/5 & 7.0 \\
\texttt{90e8f4b} & MoE: 8 experts + LLM router & 23/30 & --- \\
\texttt{2cf655d} & Evaluation scene-awareness & 5/5 & 7.4 \\
\texttt{77c6ea0} & Two-wave dispatch + rework & 4/5 & 7.6 \\
\bottomrule
\end{tabular}}
\end{table}

The key architectural decisions were: (1)~replacing the central solver with LLM-orchestrated experts (commit \texttt{107d75d}, 4/5 pass, 7.2 overall); (2)~introducing scene-aware expert dispatch (commit \texttt{7edcb97}, maintained 4/5 pass); (3)~replacing rule-based scene classification with LLM-based MoE routing (commit \texttt{90e8f4b}, improved 30-scenario pass rate from 18\% to 77\%); and (4)~making evaluation scene-aware so that food routes are not penalized for lacking scenic viewpoints (commit \texttt{2cf655d}, first 5/5 pass on core scenarios). The two-wave dispatch architecture (commit \texttt{77c6ea0}) improved the 5-scenario score from 67.4 to 76.3 but initially reduced pass rate from 3/5 to 4/5 before further optimization resolved it.

\subsection{Optimization Journey}
\label{sec:journey}

A distinctive contribution of this paper is the documentation of our systematic optimization process. Starting from a baseline score of 65.2/100 on 100 scenarios with 8 F-grade scenarios, we achieved 81.2/100 with 0 F-grade scenarios through seven rounds of targeted fixes. The final 100-scenario evaluation (2026-06-02) confirms 90\% pass rate with an average score of 74.3/100 and grade distribution S:12, A:24, B:23, C:30, D:10, F:1. Table~\ref{tab:progression} shows the score progression during development.

\begin{table}[!t]
\centering
\caption{Score progression across optimization rounds. All scores from 100-scenario stress tests.}
\label{tab:progression}
\resizebox{\columnwidth}{!}{\begin{tabular}{@{}llcl@{}}
\toprule
\textbf{Round} & \textbf{Fix Category} & \textbf{Score} & \textbf{Key Insight} \\
\midrule
0 & Baseline & 65.2 & 8 F-grade scenarios \\
1 & Evaluation bug (fake zeros) & 65.2 & Score table had wrong columns \\
2 & LLM spatial blindness & 70.1 & Add coordinates to eval prompt \\
3 & Sub-category misclassification & 71.8 & ``Herbal tea'' $\neq$ ``tea restaurant'' \\
4 & Prompt-based refusal & 76.5 & ``You have the right to refuse'' \\
5 & Time gap correction & 78.9 & \texttt{\_smooth\_times()} post-processing \\
6 & Route inflation prevention & 80.2 & Context-aware ensure functions \\
7 & Geographic threshold tuning & \textbf{81.2} & Stricter thresholds for destination \\
\bottomrule
\end{tabular}}
\end{table}

\subsubsection{Round 1: Evaluation Bug (Fake Zeros)}

\textbf{Symptom.} Five evaluation dimensions showed near-zero scores for food, destination, and speed-run scenarios.

\textbf{Root cause.} The LLM evaluation prompt contained 4 dimensions while the rule-based evaluation used 6 dimensions. When merging results, missing dimensions were filled with zeros. Commit \texttt{25bf6a0} unified the LLM evaluation prompt to 6 standard dimensions with a \texttt{\_normalize\_dims()} bridging function.

\textbf{Lesson.} When a dimension shows zero, suspect data collection before suspecting algorithm logic. This finding aligns with recent work on LLM-as-judge reliability~\cite{gu2024llmjudge}, which identifies evaluation methodology errors as a common source of misleading results.

\subsubsection{Round 2: LLM Spatial Blindness}

\textbf{Symptom.} Food routes received geographic continuity scores of 45/100, judged as ``severe geographic jumping.'' Manual verification showed all stops were within 2km of each other.

\textbf{Root cause.} The evaluation prompt provided POI names and categories but \textbf{no coordinates}. The LLM guessed distances from names, then confidently judged incorrectly. This is consistent with findings on spatial hallucination in LLMs~\cite{zhang2025spatial}.

\textbf{Fix.} Added coordinates and inter-stop distances to the evaluation prompt, with hard rules: $<$3km = 90--100, 3--8km = 70--89, $>$15km = 0--49. Commit \texttt{487f99a}.

\textbf{Lesson.} LLMs fabricate facts they don't have. Provide ground truth (coordinates), and they judge correctly.

\subsubsection{Round 3: Sub-Category Misclassification}

\textbf{Symptom.} ``De Yi Ji Herbal Tea'' was classified under ``tea restaurant/dessert'' sub-category, causing apparent duplication with actual tea restaurants.

\textbf{Root cause.} Sub-category matching used keyword containment: ``herbal tea'' contains ``tea'', which matched the tea restaurant category.

\textbf{Fix.} Added priority matching for herbal tea keywords before general tea matching. Commit \texttt{18792db}.

\textbf{Lesson.} Chinese keyword matching requires layered precision: check exact terms before fuzzy containment.

\subsubsection{Round 4: Prompt-Based Refusal}

\textbf{Symptom.} ``Zhuhai island-hopping in one day (3 islands)'' produced a route spanning three islands, despite ferry transit alone requiring 6 hours.

\textbf{Initial approach.} We considered adding a dedicated \textit{feasibility gate} architectural component.

\textbf{Actual fix.} Six lines of prompt instruction at the top of the synthesis prompt:

\begin{quote}
\textit{``You have the right to refuse unreasonable requests: If total time $<$ 3 hours, schedule at most 2--3 stops. If islands are separated by $>$2 hours ferry transit, keep only one island. Better a short, excellent route than a long, terrible one.''}
\end{quote}

Destination-focused scenarios improved from B~(78) to S~(92); speed-run scenarios from F~(42) to A~(85). Commit \texttt{487f99a}.

\textbf{Lesson.} Before adding architecture, try prompt engineering. This aligns with the ``prompt before architecture'' principle~\cite{white2023prompt}: LLMs can learn to refuse if explicitly instructed.

\subsubsection{Round 5: Time Allocation Gaps}

\textbf{Symptom.} Speed-run routes showed 2-hour gaps between consecutive stops (e.g., stop~8 departs 14:45, stop~9 arrives 17:00).

\textbf{Root cause.} The LLM-generated time allocations were trusted without validation. The LLM has no physical clock and doesn't know that 15 minutes cannot cover 3km.

\textbf{Fix.} Introduced \texttt{\_smooth\_times()} (Section~\ref{sec:method}) to algorithmically validate and correct LLM time proposals. Commit \texttt{487f99a}.

\textbf{Lesson.} LLM proposes, algorithm validates. Never trust LLM temporal reasoning without verification.

\subsubsection{Round 6: Route Inflation}

\textbf{Symptom.} Food routes were capped at 6 stops, but \texttt{\_ensure\_poi\_in\_route} added non-food POIs back, inflating to 9 stops with irrelevant entries.

\textbf{Root cause.} Four sequential \texttt{ensure} functions each independently added POIs without awareness of other functions' additions or scenario-specific constraints.

\textbf{Fix.} Made \texttt{\_ensure\_poi\_in\_route} scenario-aware: skip for food scenarios. Added station count upper bound. Commit \texttt{5721bc7}.

\textbf{Lesson.} Ensure functions are ``well-intentioned troublemakers.'' Each makes sense in isolation; combined, they inflate. Give each ensure function awareness of its boundaries.

\subsubsection{Round 7: Geographic Threshold Mismatch}

\textbf{Symptom.} Destination-focused route: Chimelong Ocean Kingdom $\rightarrow$ restaurant 12km away. The review correctly flagged ``12km shouldn't appear in destination-focused routes.''

\textbf{Root cause.} The \texttt{\_geo\_reroute} threshold was 15km for all scenarios. 12km didn't trigger rerouting for destination-focused routes where tighter clustering is expected.

\textbf{Fix.} Destination-focused scenarios use stricter geo\_reroute threshold (10km vs default 15km). Commit \texttt{0a44c47}.

\subsection{Per-Scenario Score Changes}

\begin{table}[!t]
\centering
\caption{Per-scenario scores before and after optimization.}
\label{tab:scenario_changes}
\resizebox{\columnwidth}{!}{\begin{tabular}{@{}lccc@{}}
\toprule
\textbf{Scenario} & \textbf{Initial} & \textbf{Final} & \textbf{Key Fix} \\
\midrule
Food Exploration & D (58) & B (72) & Sub-type diversity + ensure boundaries \\
Destination (Chimelong) & B (78) & S (92) & Stricter geo threshold + refusal \\
Speed-Run Touring & F (42) & A (85) & Time smoothing + stop capping \\
Leisure & B (76) & B (72) & Stable \\
Sightseeing & B (72) & A (85) & Coordinate-aware evaluation \\
\bottomrule
\end{tabular}}
\end{table}

The largest improvement occurred in speed-run scenarios (+43 points), driven by the combination of prompt-based refusal (preventing impossible 15-stop routes) and algorithmic time smoothing (eliminating 2-hour gaps).

\subsection{Geographic Continuity Improvement}

\begin{table}[!t]
\centering
\caption{Geographic continuity scores before and after adding coordinates to evaluation.}
\label{tab:geo_improvement}
\resizebox{\columnwidth}{!}{\begin{tabular}{@{}lcc@{}}
\toprule
\textbf{Scenario} & \textbf{Before (no coords)} & \textbf{After (with coords)} \\
\midrule
Food Exploration & 45 & \textbf{100} \\
Leisure & 70 & \textbf{88} \\
Sightseeing & 60 & \textbf{95} \\
\bottomrule
\end{tabular}}
\end{table}

Table~\ref{tab:geo_improvement} shows the dramatic improvement in geographic continuity scores (e.g., 45$\rightarrow$100 for food) demonstrates that the LLM was not evaluating route quality but rather guessing distances from POI names. Providing ground-truth coordinates resolved this entirely.

\subsection{Lessons from Failed Optimizations}

Beyond the seven successful rounds, we attempted 14 additional optimizations that failed or provided marginal gains. We document these because negative results are as instructive as positive ones. Table~\ref{tab:failed_opts} summarizes the key failures.

\begin{table}[!t]
\centering
\caption{Failed or marginal optimization attempts. Data from optimization log (commit \texttt{d323625}).}
\label{tab:failed_opts}
\resizebox{\columnwidth}{!}{\begin{tabular}{@{}llll@{}}
\toprule
\textbf{Attempt} & \textbf{Direction} & \textbf{Result} & \textbf{Root Cause} \\
\midrule
\#3 & Expert rule-ification & 3/5, $-$0.4 & Wrong baseline assumption \\
\#5 & Short prompt & $-$0.2 & Lost context information \\
\#7 & Diversity post-processing & 3/5, $-$0.4 & Data gap, not pipeline \\
\#12 & Diversity prompt tuning & $\pm$0 & LLM variance ate gains \\
\#14 & Cap with food preservation & 0/5, $-$1.0 & Broke geographic ordering \\
\bottomrule
\end{tabular}}
\end{table}

Three broader lessons emerge:

\textbf{Lesson 1: Verify baseline configuration before optimizing.} Attempt~\#3 replaced LLM experts with rule-based logic, assuming the baseline used expensive DeepSeek models. In fact, the baseline already used inexpensive Qwen models (commit \texttt{477531a}). The ``optimization'' replaced a cheap intelligent decision with a rigid rule.

\textbf{Lesson 2: Data gaps masquerade as pipeline failures.} Attempts~\#7 and~\#12 tried to improve scene diversity through post-processing and prompt tuning. The actual bottleneck was the POI database: entertainment POIs comprised only 3.5\% of the database, and natural-scenery POIs were virtually absent (0.05\%). After supplementing 52 food POIs (commit \texttt{e2621aa}), food-scenario scores improved from 50.5 to 57.6 without any pipeline changes.

\textbf{Lesson 3: LLM variance obscures small improvements.} The same code and prompt produce $\pm$1 point variation across runs. This means any single-run improvement of $<$1 point cannot be reliably attributed to the change. We addressed this by running 3 evaluations and reporting the median, but acknowledge that larger evaluation sets would strengthen confidence.

\subsection{Data-Driven Improvement}

A recurring finding is that data quality often matters more than algorithmic sophistication. Table~\ref{tab:data_improvements} shows the impact of POI database enrichment.

\begin{table}[!t]
\centering
\caption{Impact of POI database enrichment on evaluation scores.}
\label{tab:data_improvements}
\resizebox{\columnwidth}{!}{\begin{tabular}{@{}llcc@{}}
\toprule
\textbf{Enrichment} & \textbf{Commit} & \textbf{Before} & \textbf{After} \\
\midrule
52 food POIs (7 categories) & \texttt{e2621aa} & 50.5 & 57.6 \\
42 entertainment/nature POIs & \texttt{7e7da08} & 6.4 & 6.8 \\
30 remote-area food POIs & \texttt{b9c7079} & --- & --- \\
Category diversity enforcement & \texttt{c00f63d} & 93.5 & 99.3 \\
\bottomrule
\end{tabular}}
\end{table}

The category diversity enforcement (commit \texttt{c00f63d}) is particularly instructive: by simply ensuring that routes include at least one restaurant POI, the rule-based score jumped from 93.5 to 99.3 and the LLM-judged score from 7.0 to 8.0---a purely data-level fix with no algorithmic changes.

\subsection{Real-Time Performance}

Table~\ref{tab:scenario_latency} reports per-scenario latency from the 2026-06-02 test run.

\begin{table}[!t]
\centering
\caption{Per-scenario latency (median of 2 runs, 2026-06-02).}
\label{tab:scenario_latency}
\resizebox{\columnwidth}{!}{\begin{tabular}{@{}lc@{}}
\toprule
\textbf{Scenario} & \textbf{Elapsed (s)} \\
\midrule
S1: Couple Tour & 71.3 \\
S2: Ocean Kingdom & 38.9 \\
S3: Food Tour & 70.4 \\
S4: Speed-Run & 39.0 \\
S5: Leisure & 52.2 \\
\midrule
\textbf{Average} & \textbf{54.4} \\
\bottomrule
\end{tabular}}
\end{table}

The high variance (38.9--71.3s) reflects LLM API latency rather than algorithmic complexity. Scenarios requiring more expert proposals (S1, S3) take longer due to additional LLM calls. The average end-to-end latency of 54.4s is dominated by 7--12 sequential LLM API calls (see Table~\ref{tab:model_compare}). Using the Qwen3.5-35B-A3B (iFlytek Spark) model, the estimated per-request cost is \$0.001, compared to \$0.006 with DeepSeek-V4-Flash (commit \texttt{a576cc5}: average 48.2s with DeepSeek vs.\ 54.4s with Qwen3.5-35B-A3B).

\subsection{Evaluation Reliability}

To assess the reliability of our LLM judge, we compare two independent evaluation runs of the same 5 core scenarios (2026-06-02). Table~\ref{tab:consistency} reports the score differences between runs.

\begin{table}[!t]
\centering
\caption{LLM judge reliability: score difference between two independent evaluation runs.}
\label{tab:consistency}
\resizebox{\columnwidth}{!}{\begin{tabular}{@{}lcccc@{}}
\toprule
\textbf{Scenario} & \textbf{$\Delta$Intent} & \textbf{$\Delta$POI} & \textbf{$\Delta$Geo} & \textbf{$\Delta$Overall} \\
\midrule
S1: Couple Tour & 1 & 1 & 3 & 1 \\
S2: Ocean Kingdom & 0 & 0 & 0 & 0 \\
S3: Food Tour & 0 & 1 & 1 & 0 \\
S4: Speed-Run & 0 & 1 & 2 & 0 \\
S5: Leisure & 1 & 0 & 2 & 1 \\
\midrule
\textbf{Max $\Delta$} & \textbf{1} & \textbf{1} & \textbf{3} & \textbf{1} \\
\bottomrule
\end{tabular}}
\end{table}

The overall score shows $\leq$1 point variation across runs, confirming that our scenario-aware scoring rubrics produce consistent evaluations. Geographic continuity shows the largest variation ($\leq$3 points), consistent with the optimization log observation that LLM variance is the dominant noise source (commit \texttt{d323625}). This justifies our conservative approach of reporting the lower-score run.

\subsection{Baseline Comparison}
\label{sec:baseline_comparison}

To contextualize CityFlow's performance, Table~\ref{tab:baseline_compare} compares against three baselines on the same 100 test scenarios: (1)~\textbf{Greedy Nearest-Neighbor}, which selects the top-$k$ POIs by category match score and sorts them by geographic proximity, (2)~\textbf{Single-LLM}, which generates a complete route in a single LLM call without multi-agent decomposition, and (3)~\textbf{Single-LLM+RAG}, which augments the single LLM call with retrieved real POI context from the database to ground its recommendations.

\begin{table}[!t]
\centering
\caption{Baseline comparison on 100 test scenarios. ``Hallucination'' refers to POIs not present in the database. ``Diversity'' counts unique routes per scenario type.}
\label{tab:baseline_compare}
\resizebox{\columnwidth}{!}{\begin{tabular}{@{}lccccc@{}}
\toprule
\textbf{Metric} & \textbf{Greedy} & \textbf{Single-LLM} & \textbf{Single-LLM+RAG} & \textbf{CityFlow} \\
\midrule
Pass Rate ($\geq$3 stops) & 100\% & 100\% & 99\% & 90\% \\
Avg Route Score & 100.0 & 66.9 & 83.8 & 74.3$^{*}$ \\
Hallucination Rate & 0.0\% & 49.2\% & 0.0\% & 0.0\% \\
Route Diversity & 1/20 per type & 20/20 per type & 15--18/20 & 20/20 per type \\
Avg Stops per Route & 8.0 & 5.3 & 5.0 & 7.2 \\
\bottomrule
\end{tabular}}
\vspace{2pt}
{\footnotesize $^{*}$CityFlow score from LLM-as-Judge (7-dimension, 0--100); other scores from simplified rule-based rubric. Scores are not directly comparable across evaluation methods.}
\end{table}

\textbf{Greedy Nearest-Neighbor} achieves a perfect simplified score (100.0) and zero hallucinations by selecting exclusively from the real POI database. However, it produces \textbf{identical routes for all scenarios of the same type}---every food scenario returns the same 8 restaurants, and every sightseeing scenario returns the same 8 scenic spots. This complete lack of personalization makes it impractical for real-world use despite its high score. The 100.0 score also reveals a limitation of the simplified scoring rubric: it rewards POI quality and category matching but does not penalize route homogeneity.

\textbf{Single-LLM} achieves 100\% pass rate with full route diversity (each scenario produces a unique route), but suffers from a \textbf{49.2\% hallucination rate}---nearly half of the recommended POIs do not exist in the Zhuhai POI database. Examples include ``Yijian Seafood Restaurant'', ``Hengqin Oyster King'', ``Beishan Ancient Village'', and ``Riyuebei Opera House''---plausible-sounding names that are fabrications. The average score of 66.9 reflects penalties for geographic incoherence and category mismatch when hallucinated POIs are mixed with real ones. This confirms that a single LLM lacks the domain-specific grounding needed for reliable POI recommendation.

\textbf{Single-LLM+RAG} addresses the hallucination problem by injecting retrieved POI context into the LLM prompt before generation. The result is striking: hallucination drops from 49.2\% to \textbf{0.0\%} while the average score rises from 66.9 to \textbf{83.8}. RAG achieves the highest simplified score among all baselines, validating that domain-grounded context is critical for POI recommendation. However, route diversity decreases slightly (15--18/20 vs.\ 20/20 for Single-LLM), suggesting that RAG context narrows the LLM's focus toward the most relevant retrieved POIs. This tradeoff between accuracy and diversity mirrors a fundamental tension in retrieval-augmented generation.

\textbf{CityFlow} occupies the favorable middle ground: zero hallucinations (all POIs drawn from the database), full route diversity (personalized per-scenario routes), and the highest practical quality as judged by the LLM-as-Judge evaluation (74.3). The 90\% pass rate (vs.\ 100\% for baselines) reflects CityFlow's conservative refusal of infeasible requests (e.g., ``visit 3 islands in one day''), which we argue is a feature rather than a bug---a real tourism system should refuse unreasonable plans rather than generate poor ones.

These results validate our core hypothesis: \textbf{multi-agent decomposition with expert specialization provides a complementary advantage beyond what RAG alone can achieve}. RAG eliminates hallucination and maximizes simplified scoring, but CityFlow's multi-agent architecture delivers superior personalization (full diversity) and more nuanced quality assessment through its LLM-as-Judge evaluation pipeline.
